\centerline{\bf \Huge КТО КТО в теремочке живёт...}
\centerline{\bf \Huge }

\begin{flushright}
        \scriptsize{}
\end{flushright}

\problem{1}
\textbf{Китайская теорема об остатках для двух чисел.} Докажите, что для взаимно простых чисел $a$ и $b$ и любой пары остатков $n<a$ и $k<b$ найдется число $c$ такое, что при делении на $a$ число $c$ даёт в остатке $n$, а при делении на $b$ даёт в остатке $k$. При этом такое $c$ единственно на промежутке $0 \leqslant c<a b$.

\textit{У.} Верно ли предыдущее утверждение для не взаимно простых чисел $a$ и $b$ ?

П. Очевидно.

\problem{2}
Пусть числа $a$ и $b$ взаимно просты. Докажите, что при делении чисел от 1 до $a b$ на $a$ и на $b$ получаются все возможные пары остатков.

\problem{3}
При каких целых $n$ число $n^2+3 n+1$ делится на 55 ?

\problem{4}
Найдите наименьшее натуральное число, дающее остаток 2 при делении на 3 , остаток 3 при делении на 4 , остаток 4 при делении на 5 , остаток 5 при делении на 6 и остаток 6 при делении на 7 .

\problem{5}
\textbf{Китайская теорема об остатках (KTO).} Пусть числа $m_1, m_2, \ldots, m_n$ попарно взаимно просты. Тогда для любых целых $a_1, a_2, \ldots, a_n$ найдется целое число $x$ такое, что $x \equiv a_i\left(mod \ m_i\right)$ для всех $i$. Более того, $x$ определен однозначно с точностью до прибавления кратного $M=m_1 m_2 \ldots m_n$.

\problem{6}
\textit{Полезное следствие из КТО.} Докажите, что для любых попарно взаимно простых чисел $m_1, m_2, \ldots, m_n$ и остатков $r_1, r_2, \ldots, r_n$ по модулям $m_1, m_2, \ldots, m_n$ найдутся $n$ последовательных чисел $a+1, a+2, \ldots, a+n$ таких, что $a+i \equiv r_i\left(\bmod m_i\right)$.

\problem{7}
Докажите, что найдутся 2020 последовательных натуральных чисел, каждое из которых имеет по меньшей мере три различных простых делителя.

\problem{8}
Докажите, что числа натурального ряда можно переставить местами (то есть после перестановки всё ещё каждое число должно встречаться в ряду ровно один раз) так, чтобы для всех $n$ сумма $n$ первых чисел делилась на $n$.